% Figure: straight-line vehicle performance. The tractive force a car
% can apply is the lesser of the traction limit (a flat ceiling set by
% tyre grip) and the power limit (a hyperbola F = P/v that falls with
% speed). The crossover speed separates the traction-limited launch
% from the power-limited top end.
\begin{figure}[htbp]
\centering
\begin{tikzpicture}
\begin{axis}[
  width=12cm, height=7cm,
  xlabel={speed $v$ (\si{\meter\per\second})},
  ylabel={available tractive force (\si{\kilo\newton})},
  xmin=0, xmax=70, ymin=0, ymax=14,
  axis lines=left,
  domain=5:70, samples=140,
  legend pos=north east,
  legend cell align=left,
  grid=major, grid style={enggray!20},
  tick label style={font=\scriptsize},
  label style={font=\small},
  legend style={font=\scriptsize},
]
% Traction limit: mu m g = 0.9 * 1400 * 9.81 = 12.36 kN, flat.
\addplot[engred, very thick, domain=0:70] {12.36};
\addlegendentry{traction limit $\mu m g$}
% Power limit: P/v, P = 150 kW = 150000 W, in kN -> 150/v.
\addplot[engblue, very thick] {150/x};
\addlegendentry{power limit $P/v$}
% Resistance (drag + rolling): R = 0.2 + 0.0006*v^2 kN (rough).
\addplot[enggray, thick, dashed, domain=0:70] {0.2 + 0.0006*x*x};
\addlegendentry{road load (drag + rolling)}
% crossover of traction and power limits: 150/v = 12.36 -> v = 12.1 m/s
\addplot[only marks, mark=*, engblack, mark size=2pt] coordinates {(12.1,12.36)};
\node[font=\scriptsize, anchor=south west] at (axis cs:12.1,12.36)
  {crossover $v\approx 12\,\si{\meter\per\second}$};
% top speed: power limit meets road load near v = 61 m/s
\addplot[only marks, mark=square*, engblack, mark size=2pt] coordinates {(61,2.46)};
\node[font=\scriptsize, anchor=south east] at (axis cs:61,2.46)
  {top speed};
\end{axis}
\end{tikzpicture}
\caption[Traction-limited and power-limited vehicle performance]{The
straight-line performance of a car is governed by two ceilings on the
tractive force. Below the crossover speed the tyres set the limit: the
ground can return at most $\mu m g$ regardless of engine output, so the
launch is traction-limited and the acceleration is the friction
maximum. Above the crossover the engine sets the limit: a fixed power
$P$ delivers force $P/v$, a hyperbola that falls as the car speeds up,
so the high-speed pull is power-limited. The two limits cross near
$12\,\si{\meter\per\second}$ for the numbers in the worked example.
Top speed is reached where the power-limited pull (blue) drops to meet
the road load (grey dashed), the sum of aerodynamic drag and rolling
resistance; no surplus force remains to accelerate.}
\label{fig:vol03ch04:traction-power}
\end{figure}
